\chapter{What's a Fair Start?}

In this eighth lecture, Michael Sandel turns directly to distributive justice: how income, wealth, power, and opportunity should be distributed, and according to what principles. These notes, prepared in the LazyLearn track and curated by LazyingArt LLC, follow the lecture's order closely. The argument unfolds in stages: first the original position and the veil of ignorance, then the two principles of justice, then the challenge from merit and effort, and finally the distinction between moral desert and entitlements to legitimate expectations. Because no validated mathematical board frames survive for this episode, every displayed equation below is a transcript-backed schematic reconstruction of Sandel's spoken argument rather than a transcription of a visible blackboard.

\section{From the Original Position to the First Principle}

Sandel opens by reminding us where we are in the larger Rawls sequence. Last time, the work was to see why principles of justice should be chosen from a hypothetical contract. This time, the question is narrower and sharper: what principles would actually be chosen there?

We can summarize the opening setup in the compact chain that governs the whole lecture:
\begin{equation}
\text{veil of ignorance}
\Longrightarrow
\text{original position of equality}
\Longrightarrow
\text{choice of principles of justice}.
\end{equation}

The point of the veil is not mystery for its own sake. It is to prevent us from tailoring principles to our present luck. We do not know whether we will be rich or poor, talented or ordinary, healthy or infirm, members of a majority or a minority. Once that device is in place, Sandel asks the first live question: what major rival would be rejected first?

The answer is utilitarianism. Rawls asks whether people in the original position would choose to govern their collective lives by the greatest good for the greatest number. Sandel's answer is no, and the reason is immediately motivational. Once the veil rises, each of us will want to live with dignity. Any of us may turn out to be vulnerable, unpopular, or exposed. No one would rationally agree to a principle that could sacrifice his or her fundamental rights for aggregate welfare.

\subsection{Question \& Answer}

\paragraph{Question.} Why would parties behind the veil reject utilitarianism?

\paragraph{Answer.} Because behind the veil we do not know whether we will later belong to an oppressed minority. A utilitarian principle may permit the happiness of the many to override the rights of the few. Sandel emphasizes that this is precisely the risk no one would knowingly run from a position of fair choice.

The structure of the rejection can be written as
\begin{equation}
\text{utilitarianism}
\Longrightarrow
\text{possible sacrifice of minorities}
\Longrightarrow
\text{rejected}.
\end{equation}

Rawls's criticism is that utilitarianism does not take seriously the distinction between persons. From that rejection Sandel derives the first principle:
\begin{equation}
P_1=\text{equal basic liberties}.
\end{equation}

The spoken list in the transcript is momentarily garbled, but its stable content is clear: freedom of speech, freedom of assembly, religious liberty, freedom of conscience, and the like. The order matters. We do not first maximize income and then see what liberties remain. Equal basic liberties come first, and only after that do we turn to social and economic inequalities.

\section{From Strict Equality to the Difference Principle}

With the first principle secured, Sandel shifts the lecture from liberty to distribution. What should we agree to concerning income and wealth once we do not know where we will land? Here he deliberately stages an initial answer and then replaces it.

The first answer is strict equality. If we are uncertain about our future place in society, equal shares seem like the safest rule:
\begin{align}
\text{strict equality} &= \text{equal distribution of income and wealth},\\
\text{strict equality} &\not\Longrightarrow \text{best outcome for the least well-off}.
\end{align}

This second line is the crucial turn. Equal shares may be safe, but they may not be best. Even someone who winds up at the bottom might do better under a qualified inequality than under strict equality.

Rawls's second principle, as Sandel presents it, is therefore not an argument for eliminating all inequality:
\begin{equation}
P_2=\text{allow only those social and economic inequalities that benefit the least well-off}.
\end{equation}

In the lecture itself, this is phrased in prose. If we want a compact shorthand for the same idea, we can write
\begin{equation}
\Delta(\text{income, wealth})>0
\qquad\text{is acceptable only if}\qquad
\Delta(\text{position of the least well-off})>0.
\end{equation}

This notation is only a careful compression of Sandel's spoken test. It is not meant as Rawls's own formal notation. What matters is the direction of the judgment: we evaluate inequality from the standpoint of those at the bottom.

Sandel immediately stress-tests the point with famous examples. Michael Jordan's salary and Bill Gates's fortune are not introduced merely as symbols of excess. They are introduced to force the institutional question: if inequalities of this magnitude are ever justified, what kind of social scheme could justify them?

\subsection{Question \& Answer}

\paragraph{Question.} Why would rational choosers accept some inequality rather than insist on equal shares?

\paragraph{Answer.} Because once we ask not merely what is safest, but what leaves the least well-off best off, equal shares are no longer automatically preferred. If a qualified inequality raises the floor, then someone choosing behind the veil might rationally prefer it to equality that leaves the floor lower. Sandel therefore presents the difference principle as a qualified principle of equality, not as a defense of laissez-faire and not as a simple equality of outcome.

\section{Merit, Family Background, and the Classroom Challenge}

At this point the lecture becomes openly dialogical. Sandel asks whether anyone rejects the claim that these two principles would be chosen behind the veil. Mike does. His challenge is exact and important: why assume that people behind the veil would legislate from the standpoint of the disadvantaged? Why not choose a merit-based system that rewards effort?

Mike gives the challenge a concrete form. Harvard, he says, is not built on random selection. People work to get there. A system that rewards effort seems to give people something to strive for, and he suggests that such a system may lift the social floor better than one aimed too directly at equality.

Kate's reply shifts the ground. She asks whether we can speak of rewarding effort as though effort emerged in a vacuum. What about the educational and family advantages that shape who is in a position to display merit in the first place? Sandel strengthens her point with the empirical finding he cites:
\begin{align}
\text{students from the bottom quarter at selective colleges} &\approx 3\%,\\
\text{students from affluent families} &> 70\%.
\end{align}

Now the lecture makes a decisive turn. Sandel says that Rawls has two arguments, not one. The first is the official argument from the original position. The second is a straightforwardly moral argument: the distribution of income, wealth, and opportunity should not be based on factors for which people can claim no credit.

That turn matters. The classroom challenge does not merely interrupt the lecture. It opens the route from the contract argument to the argument from moral arbitrariness.

\subsection{Question \& Answer}

\paragraph{Question.} Why not choose a meritocratic system behind the veil of ignorance?

\paragraph{Answer.} Because meritocracy looks fair only if the conditions under which merit is displayed are already fair. Kate's point is that they are not. Economic background, cultural capital, and educational advantage enter long before the official competition begins. Sandel uses this exchange to show that the real question is not just what rational choosers would pick under uncertainty, but whether distributive shares should rest on morally arbitrary contingencies at all.

\section{The Ladder of Rival Theories}

Having widened the frame, Sandel reconstructs Rawls's ladder of rival theories of justice. The force of the ladder lies in its sequence. Each theory corrects one arbitrariness and leaves another standing.

We can write the progression in the compressed form that the lecture itself invites:
\begin{align}
\text{feudal aristocracy} &\Longrightarrow \text{life prospects fixed by accident of birth},\\
\text{formal equality of opportunity} &= \text{careers open to talents},\\
\text{formal equality}+\text{unequal starting points} &\Longrightarrow \text{unfair race},\\
\text{fair equality of opportunity} &\Longrightarrow \text{same starting line},\\
\text{same starting line}+\text{natural lottery} &\not\Longrightarrow \text{full justice}.
\end{align}

Feudal aristocracy is the easiest case. It is unjust because life prospects are fixed by birth. Formal equality of opportunity improves on that by opening careers to all. But it does not go far enough, because those who are permitted to enter the race do not begin from equal positions. Family wealth, schooling, upbringing, and social support all affect the result.

Fair equality of opportunity corrects that next arbitrariness by trying to bring everyone to the same starting line. Sandel follows Rawls in treating this as genuine moral progress. But once the runners begin together, the next question arises immediately: who wins? The fastest runners. And is it their doing that they have the natural speed to win?

This is the point at which Rawls moves beyond meritocracy. The lecture is especially careful here not to collapse into caricature. Rawls does not say that the gifted must be handicapped or that equality requires leveling all outcomes. Rather, he shifts the question from whether talents may be exercised to the terms on which their fruits may be held:
\begin{equation}
\text{no leveling equality}
\qquad
\text{but}
\qquad
\text{change the terms on which talent-based gains are held}.
\end{equation}

The guiding statement becomes
\begin{equation}
\text{people may gain from good fortune}
\Longrightarrow
\text{only on terms that help the least well-off}.
\end{equation}

\subsection{Question \& Answer}

\paragraph{Question.} Why is formal equality of opportunity, and even fair equality of opportunity, still not enough?

\paragraph{Answer.} Formal equality leaves social starting points unequal. Fair equality of opportunity corrects that, but it still leaves the natural lottery intact. If the moral complaint is really about arbitrary advantages, then equal permission to compete is not enough, and even equal starting conditions are not enough. We must also ask what claim the naturally gifted have on the benefits that flow from their luck. Rawls's answer is that they may enjoy those benefits, but only on terms that improve the condition of those who are less fortunate.

\section{Natural Lottery, Re-entry, and Incentives}

At this point the lecture visibly circles back on itself. Sandel briefly poses a dramatic salary example, recalls the earlier birth-order poll and the rival theories of distributive justice, and then re-enters the argument through the concrete pay differentials of ordinary life. That doubling back is worth preserving because it shows what he is doing pedagogically: he is not abandoning the theory, but translating it into a form in which objections can now be heard clearly.

The difference principle is tested not on small variations but on striking inequalities. Sandel asks what the average school teacher makes, and then what David Letterman makes.

\paragraph{Worked example.} Let us write Sandel's comparison as
\begin{align}
T &\approx \$42{,}000,\\
L &\approx \$31{,}000{,}000,\\
L-T &\approx \$30{,}958{,}000.
\end{align}

The Rawlsian question is not whether \(L>T\). That is obvious. The question is whether the basic structure of society is arranged in such a way that the large excess represented by \(L-T\) works, directly or indirectly, to the advantage of the least well-off.

The same structure appears in Sandel's second example:
\begin{align}
\text{Supreme Court justice salary} &< \$200{,}000,\\
\text{Judge Judy salary} &\approx \$25{,}000{,}000.
\end{align}

Again, Rawls's answer is conditional. Such inequalities are not justified simply because the market produces them. They are justified only if the institutional scheme surrounding them improves the situation of those at the bottom.

Sandel then turns to the first major objection. What about incentives? If we tax away too much of what the talented can earn, will they cease to produce? Will Michael Jordan stop playing, David Letterman stop performing, or CEOs move elsewhere?

\subsection{Question \& Answer}

\paragraph{Question.} Can the difference principle allow incentives without collapsing into laissez-faire?

\paragraph{Answer.} Yes. Tim's answer in the lecture, which Sandel endorses, is that Rawls can allow incentive-based inequalities. But the standpoint of judgment is decisive. We do not ask only whether incentives enlarge total output. We ask whether the incentive structure improves the condition of the least well-off.

That is the distinction Sandel wants us to keep firmly in view:
\begin{equation}
\text{incentives judged from standpoint of the least well-off}
\neq
\text{incentives judged by total output alone}.
\end{equation}

Sandel even gives the point a more textual Rawlsian form. The naturally advantaged are not to gain merely because they are more gifted, but only in ways that cover training, sustain useful activity, and help the less fortunate as well. Incentives, then, are not a decisive objection to the difference principle. They are already built into its correct application.

\section{Self-Ownership and the Libertarian Objection}

Once incentives have been answered, Sandel names two deeper objections. One comes from meritocratic defenders of effort and desert. The other comes from libertarians who invoke self-ownership. He takes the libertarian objection first.

Here Sandel distinguishes an easier argument from a harder one. The easier argument, associated with Milton Friedman, says that life is unfair and that any attempt to correct nature will lead to a disastrous leveling equality of outcome. Sandel answers that Rawls is not trying to equalize outcomes in that crude sense. The relevant question is not whether nature distributes talents equally, but what institutions do with natural facts.

This is the point at which Sandel quotes one of Rawls's most famous formulations:
\begin{equation}
\text{natural distribution of talents is neither just nor unjust}.
\end{equation}
And the companion thought is
\begin{equation}
\text{justice}=\text{how institutions deal with those facts}.
\end{equation}

The deeper libertarian objection, associated with self-ownership, is harder. Even if one grants that public schools and fair starting conditions are desirable, why may society tax the earnings of talented people to create them? If the state takes part of David Letterman's income against his will, does it not simply steal from him?

\subsection{Question \& Answer}

\paragraph{Question.} If talents are mine, why may society tax the gains that flow from them?

\paragraph{Answer.} Sandel's reply is careful. Rawls does not deny the importance of the person, and he does not permit the state to commandeer lives at will. The first principle remains in force: equal basic liberties are secure. What Rawls denies is the stronger thesis that because our talents are ours, we have an unlimited moral claim to every market return that flows from them. Those returns depend not only on the talents themselves, but on social institutions and on the society in which those talents happen to command value.

So the issue is not whether liberty survives. It does. The issue is whether self-ownership yields a privileged claim to the full market value of our endowments. Rawls's answer, as Sandel presents it, is no.

\section{Effort, Contribution, and Moral Desert}

The lecture's final major movement returns to the meritocratic objection, but now with the later objections in view. Sandel has already raised one version of the issue in the classroom discussion about birth order, and he now gathers it into a more explicit reply.

The first part of the reply is familiar by now. Even effort is not purely our own doing. Work ethic, ambition, conscientious striving, and discipline all arise within family and social conditions for which we can claim no decisive credit. That is why the birth-order poll mattered. Sandel uses it to dramatize the point that even the traits we are most tempted to treat as self-made are entangled with contingency.

The second part of the reply is sharper. Meritocrats say they reward effort, but when we test that claim, what they actually reward is contribution. Sandel makes the point with his construction-worker example.

\paragraph{Worked example.} Suppose two workers complete the same job:
\begin{align}
\text{Worker A: }& 4\ \text{walls in}\ 1\ \text{hour},\\
\text{Worker B: }& 4\ \text{walls in}\ 3\ \text{days}.
\end{align}

No defender of meritocracy would conclude that Worker B deserves more simply because the task cost him more exertion. What is being rewarded is not effort as such, but productive contribution. And that brings us back to talent.

So the lecture takes one further step:
\begin{equation}
\text{contribution in a market economy}
\Longrightarrow
\text{depends on what society happens to prize}.
\end{equation}

This is Sandel's second contingency, and it is one of the lecture's deepest points. Even if we granted a clean claim to our talents and our efforts, the benefits they bring still depend on social demand. David Letterman's style is valuable because he lives in a society that prizes it. In another society, the same trait might earn little or nothing. That is why Sandel can write, in effect,
\begin{equation}
\text{supply and demand}\not\Longrightarrow \text{moral desert}.
\end{equation}

The climactic distinction now comes into focus:
\begin{equation}
\text{moral desert}\neq\text{entitlements to legitimate expectations}.
\end{equation}

Sandel slows the lecture down at exactly this point and distinguishes two different games:
\begin{align}
\text{lottery} &\Longrightarrow \text{entitlement without desert},\\
\text{game of skill} &\Longrightarrow \text{entitlement plus a further desert question}.
\end{align}

If our lottery number comes up, we are entitled to the winnings under the rules, but we do not therefore morally deserve to win. If a team wins the World Series, it is entitled to the trophy, but there remains a meaningful question about whether it deserved to win. Rawls's argument, as Sandel presents it, is that distributive justice belongs to the first register rather than the second. A just scheme secures what people are entitled to expect under just institutions; it does not distribute shares in proportion to intrinsic worth.

\subsection{Question \& Answer}

\paragraph{Question.} Is distributive justice about moral desert, or only about what institutions entitle us to expect?

\paragraph{Answer.} Sandel's answer on Rawls's behalf is that distributive justice is not a matter of moral desert. It is a matter of legitimate expectation within a just institutional order. We may be entitled to what the rules assign us, but it is a conceit to infer from that entitlement that we therefore deserve, in some prior moral sense, a society that happens to reward the qualities we possess.

\subsection{Question \& Answer}

\paragraph{Question.} Are elite educational opportunities rewards for desert, or entitlements that require public justification?

\paragraph{Answer.} The lecture ends by extending the entire Rawlsian test from income and wealth to opportunities and honors. Seats at elite colleges and universities are not simply prizes for the morally deserving. They are also social positions and opportunities, and their justification depends on whether the institutional scheme that distributes them can be defended from the standpoint of those who are least advantaged.

That final move can be written as
\begin{equation}
\text{claims to elite educational opportunity}
\Longrightarrow
\text{subject to the same Rawlsian test as claims to income}.
\end{equation}

Sandel deliberately leaves the matter open. The question of admissions is not resolved here; it is prepared. The same standard that we have applied to Michael Jordan, David Letterman, Bill Gates, and Judge Judy is now carried over to elite universities. That is the bridge to the lecture on affirmative action that follows.

\section{Summary}

The lecture unfolds as an ordered argument. We begin with the original position and the veil of ignorance, reject utilitarianism, and arrive at the first principle of equal basic liberties. We then move from strict equality to the difference principle, which allows inequalities only when they improve the condition of the least well-off. That already gives us a powerful way to think about income and wealth, but the lecture does not stop there.

Mike's objection and Kate's reply widen the frame from the official contract argument to the argument from moral arbitrariness. Rawls's ladder then carries us from feudal aristocracy to formal equality of opportunity, to fair equality of opportunity, and finally beyond meritocracy to the democratic conception defined by the difference principle. The later objections sharpen the point rather than undo it. Incentives matter, but only from the standpoint of those at the bottom. Natural talents are neither just nor unjust; what matters is the structure of institutions. Effort does not rescue moral desert, because even effort is contingent and because market reward tracks contribution, which itself depends on what society happens to prize.

So the lecture closes on Rawls's hardest distinction: not moral desert, but entitlements to legitimate expectations. And once that distinction is clear, the chapter opens outward, from income and wealth to honors and opportunities, from salary to admissions, and from distributive justice in general to the specific debate about affirmative action that Sandel takes up next.